\begin{abstract}
The distribution of galaxies holds more cosmological information than the standard two-point clustering statistics can capture: the non-linear cosmic web of nodes, filaments, walls and voids is also sensitive to the growth of structure, the matter content, dark energy and neutrinos. We ask how much of this extra information can be recovered with ASTRA, an algorithm that assigns an environment to every point from its neighbours in a Delaunay tessellation of the galaxies and of a random catalogue, without any smoothing scale. We split galaxies and random points by ASTRA environment and measure their auto- and cross-correlations, and compare them with the galaxy two-point function and with density-split clustering, the reference environment statistic. We build the covariances and cosmological derivatives needed for Fisher forecasts from three simulation suites: DESI-like galaxy mocks spanning 85 cosmologies (AbacusSummit), 1786 independent mocks of the official DESI programme, and the Quijote simulations. Along the way, we train a neural-network emulator of all the statistics and use it for an end-to-end inference test. ASTRA environment statistics constrain $\omega_b$, $\omega_{\rm cdm}$, $n_s$, $\sigma_8$, $w_0$ and $w_a$ about twice as well as the two-point function. They carry as much information as density split, and the two are complementary: combining them improves the constraints further. The original ASTRA cosmic-web classes, tested on Quijote, constrain $\sigma_8$ two to three times better than density split. We also describe the main difficulties of this kind of forecast (estimating derivatives from emulator suites, choosing a covariance, handling noisy derivatives) and how we dealt with them.
\end{abstract}
